% ============================================================================
% CAPÍTULO 3: MAXWELL Y LA LUZ
% ============================================================================
\chapter{Maxwell y la Luz}
\label{cap:maxwell}

\begin{center}
\textit{``La velocidad del universo''}
\end{center}

\vspace{0.5cm}

\begin{tcolorbox}[colback=white, colframe=kleinblue, boxrule=2pt]
\centering
\textit{``Fue así como la luz fue entendida como una onda electromagnética.''}\\
--- James Clerk Maxwell, 1865\\[0.3cm]
\textit{``Y la velocidad de esa onda está escrita en la topología del cosmos.''}\\
--- Teoría Klein, 2026
\end{tcolorbox}


% ============================================================================
\section{Las Ecuaciones que Unificaron Electricidad y Magnetismo}
% ============================================================================

\begin{analogiabox}
\textbf{El Matrimonio de Dos Fuerzas}

Antes de Maxwell, la electricidad y el magnetismo eran como dos
personas que vivían en la misma casa pero nunca se hablaban.

\begin{itemize}
    \item \textbf{Electricidad:} cargas, chispas, relámpagos
    \item \textbf{Magnetismo:} imanes, brújulas, auroras
\end{itemize}

Maxwell mostró que son la MISMA cosa, como dos caras de una moneda.
Cuando una carga se mueve, crea magnetismo.
Cuando el magnetismo cambia, crea electricidad.

Y juntos, crean algo mágico: LUZ.
\end{analogiabox}

\textbf{LAS CUATRO ECUACIONES DE MAXWELL:}

\begin{nivelmatematico}
\textbf{Las Ecuaciones de Maxwell}

\begin{enumerate}
    \item $\nabla \cdot \mathbf{E} = \rho/\varepsilon_0$ \hfill Ley de Gauss (eléctrica)\\
    \textit{``Las cargas crean campos eléctricos''}

    \item $\nabla \cdot \mathbf{B} = 0$ \hfill Ley de Gauss (magnética)\\
    \textit{``No existen monopolos magnéticos''}

    \item $\nabla \times \mathbf{E} = -\partial\mathbf{B}/\partial t$ \hfill Ley de Faraday\\
    \textit{``El magnetismo cambiante crea electricidad''}

    \item $\nabla \times \mathbf{B} = \mu_0\mathbf{J} + \mu_0\varepsilon_0\partial\mathbf{E}/\partial t$ \hfill Ley de Ampère-Maxwell\\
    \textit{``La corriente y la electricidad cambiante crean magnetismo''}
\end{enumerate}

donde:
\begin{itemize}
    \item $\mathbf{E}$ = campo eléctrico
    \item $\mathbf{B}$ = campo magnético
    \item $\rho$ = densidad de carga
    \item $\mathbf{J}$ = densidad de corriente
    \item $\varepsilon_0$ = permitividad del vacío
    \item $\mu_0$ = permeabilidad del vacío
\end{itemize}
\end{nivelmatematico}

\textbf{LA REVELACIÓN DE MAXWELL:}

De estas cuatro ecuaciones, Maxwell dedujo algo extraordinario.
En el vacío (sin cargas ni corrientes), las ecuaciones predicen
una ONDA que viaja a velocidad:

\begin{equation}
c = \frac{1}{\sqrt{\varepsilon_0 \mu_0}}
\end{equation}

Cuando Maxwell calculó este valor... ¡era igual a la velocidad de la luz!

Este fue el momento ``eureka'' de la física del siglo XIX:
\textbf{La luz ES una onda electromagnética.}


% ============================================================================
\section{La Velocidad de la Luz: $c = (3 - 1/(7\pi)^2) \times 10^8$}
% ============================================================================

Ahora viene nuestra contribución. ¿Por qué $c$ tiene ese valor específico?

\begin{nivelconceptual}
\textbf{La Fórmula Klein para $c$}

\begin{equation}
\boxed{c = \left(3 - \frac{1}{(7\pi)^2}\right) \times 10^8 \text{ m/s}}
\end{equation}

\textbf{Desglose:}
\begin{itemize}
    \item $3$ = número de dimensiones espaciales
    \item $1/(7\pi)^2 = 0.00206852$ = corrección Klein
    \item $10^8$ = factor de escala (unidades SI)
\end{itemize}

\textbf{VERIFICACIÓN:}
\begin{align}
\text{Predicción:} & \quad 299,792,459 \text{ m/s}\\
\text{Observado:} & \quad 299,792,458 \text{ m/s}\\
\text{Error:} & \quad 0.0003\%
\end{align}
\end{nivelconceptual}

\begin{analogiabox}
\textbf{El Límite de Velocidad Cósmico}

Imagina una autopista con 3 carriles.
El límite de velocidad ``ideal'' sería 3 unidades.

Pero hay pequeños baches (la topología Klein) que reducen
ligeramente la velocidad máxima:

\begin{equation}
\text{Velocidad máxima} = 3 - (\text{baches}) = 3 - \frac{1}{(7\pi)^2} \approx 2.998
\end{equation}

Los ``baches'' son la corrección de 0.2\% debido a las 7 capas Klein.
\end{analogiabox}

\textbf{¿POR QUÉ $3 - 1/(7\pi)^2$ Y NO SIMPLEMENTE 3?}

El número 3 representa las dimensiones espaciales por las que puede
propagarse la luz. Pero la topología Klein introduce una pequeña
``resistencia'' o ``impedancia'' a la propagación.

Esta resistencia es $(7\pi)^{-2}$ porque:
\begin{itemize}
    \item La luz debe ``atravesar'' la estructura Klein
    \item El exponente 2 indica que es un efecto de segundo orden
    \item $7\pi$ es el factor de supresión fundamental
\end{itemize}

\textbf{IMPLICACIÓN PROFUNDA:}

Si el universo NO tuviera topología Klein, la velocidad de la luz sería
EXACTAMENTE $3 \times 10^8$ m/s. La pequeña diferencia (207,542 m/s) es la
``huella digital'' de la estructura topológica del cosmos.


% ============================================================================
\section{La Constante de Estructura Fina: $1/\alpha = 7^2\pi - 7 - \pi^2$}
% ============================================================================

La constante de estructura fina $\alpha \approx 1/137$ es quizás el número más
misterioso de toda la física.

\begin{nivelconceptual}
\textbf{¿Qué es $\alpha$?}

\begin{equation}
\alpha = \frac{e^2}{4\pi\varepsilon_0\hbar c} \approx \frac{1}{137}
\end{equation}

$\alpha$ determina:
\begin{itemize}
    \item La fuerza de la interacción electromagnética
    \item El tamaño de los átomos
    \item Los niveles de energía atómicos
    \item Si la química es posible
\end{itemize}

Si $\alpha$ fuera ligeramente diferente:
\begin{itemize}
    \item $\alpha > 1/137$: los electrones caerían al núcleo, no habría átomos
    \item $\alpha < 1/137$: los átomos serían demasiado débiles, no habría química
\end{itemize}

El valor EXACTO de $\alpha$ es crucial para nuestra existencia.
\end{nivelconceptual}

\begin{nivelmatematico}
\textbf{La Fórmula Klein para $\alpha$}

\begin{equation}
\boxed{\frac{1}{\alpha} = 7^2\pi - 7 - \pi^2}
\end{equation}

Forma alternativa:
\begin{equation}
\frac{1}{\alpha} = 7(7\pi - 1) - \pi^2
\end{equation}

\textbf{Desglose:}
\begin{itemize}
    \item $7^2\pi = 49\pi \approx 153.94$: término principal
    \item $-7$: corrección por una capa de referencia
    \item $-\pi^2 \approx -9.87$: corrección geométrica
    \item Total $\approx 137.07$
\end{itemize}

\textbf{VERIFICACIÓN:}
\begin{align}
\text{Predicción:} & \quad 137.0683\\
\text{Observado:} & \quad 137.0360\\
\text{Error:} & \quad 0.024\%
\end{align}
\end{nivelmatematico}

\begin{analogiabox}
\textbf{El Dial del Universo}

Imagina un dial con marcas de 1 a 200.

Para que existan átomos, química y vida, el dial debe estar
en EXACTAMENTE 137. Un poco más o menos, y no estarías leyendo esto.

¿Quién puso el dial en 137?

\textbf{Respuesta Klein:} NADIE. El valor 137 surge automáticamente de
$7^2\pi - 7 - \pi^2$. Es una consecuencia matemática de la topología.

No hay dial. No hay quien lo ajuste. Es geometría pura.
\end{analogiabox}

\textbf{INTERPRETACIÓN DE LA FÓRMULA:}

\begin{align}
\frac{1}{\alpha} &= 7^2\pi - 7 - \pi^2\\
&= 7 \times 7 \times \pi - 7 - \pi^2\\
&= 7 \times (7\pi - 1) - \pi^2
\end{align}

El primer término $7^2\pi$ representa la interacción ``desnuda'' entre
7 capas y 7 capas de Klein, modulada por la geometría $\pi$.

El término $-7$ resta una capa (la capa de ``referencia'').

El término $-\pi^2$ es una corrección geométrica de segundo orden.

Resultado: 137.068... vs 137.036... observado.


% ============================================================================
\section{La Impedancia del Vacío: $Z_0 = 120\pi$}
% ============================================================================

\begin{nivelconceptual}
\textbf{¿Qué es la impedancia del vacío?}

\begin{equation}
Z_0 = \sqrt{\frac{\mu_0}{\varepsilon_0}} \approx 377 \text{ }\Omega
\end{equation}

Es la ``resistencia'' que el vacío ofrece a las ondas electromagnéticas.
Determina:
\begin{itemize}
    \item La relación entre campo eléctrico y magnético en una onda
    \item Cómo las antenas irradian energía
    \item La reflexión de ondas en interfaces
\end{itemize}
\end{nivelconceptual}

\begin{nivelmatematico}
\textbf{La Fórmula Klein para $Z_0$}

\begin{equation}
\boxed{Z_0 \approx 120\pi = 5! \times \pi}
\end{equation}

donde $5! = 5 \times 4 \times 3 \times 2 \times 1 = 120$

\textbf{VERIFICACIÓN:}
\begin{align}
\text{Predicción:} & \quad 376.9911 \text{ }\Omega\\
\text{Observado:} & \quad 376.7300 \text{ }\Omega\\
\text{Error:} & \quad 0.069\%
\end{align}

\textbf{CONEXIÓN KLEIN:}

$5!$ = permutaciones de 5 elementos\\
$5$ = dimensiones de Kaluza-Klein (4 espacio-tiempo + 1 extra)
\end{nivelmatematico}

\textbf{¿POR QUÉ $5!$?}

En teoría de Kaluza-Klein, el universo tiene 5 dimensiones:
4 de espacio-tiempo + 1 dimensión extra compacta.

Las $5! = 120$ permutaciones de estas 5 dimensiones, multiplicadas
por $\pi$ (geometría), dan la impedancia del vacío.

Esto conecta la estructura dimensional del universo con sus
propiedades electromagnéticas.


% ============================================================================
\section{De Klein a Maxwell: La Derivación}
% ============================================================================

Ahora mostramos cómo las ecuaciones de Maxwell EMERGEN de Klein.

\begin{nivelmatematico}
\textbf{Derivación}

\textbf{PASO 1: La velocidad de la luz}
\begin{equation}
c = \left(3 - \frac{1}{(7\pi)^2}\right) \times 10^8 \text{ m/s} \quad \text{(fórmula Klein)}
\end{equation}

\textbf{PASO 2: La permeabilidad del vacío}
\begin{equation}
\mu_0 = 4\pi \times 10^{-7} \text{ H/m} \quad \text{(definición SI histórica)}
\end{equation}

El factor $4\pi$ es puramente geométrico (área de esfera unitaria).

\textbf{PASO 3: La permitividad del vacío}
\begin{equation}
\varepsilon_0 = \frac{1}{\mu_0 c^2} = \frac{1}{4\pi \times 10^{-7} \times \left(3 - \frac{1}{(7\pi)^2}\right)^2 \times 10^{16}}
\end{equation}

\textbf{PASO 4: Las ecuaciones de Maxwell}

Con $\varepsilon_0$ y $\mu_0$ determinados por Klein, las ecuaciones de Maxwell
quedan completamente especificadas.

\textbf{CONCLUSIÓN:}

Las ecuaciones de Maxwell son consecuencia de:
\begin{itemize}
    \item Geometría ($4\pi$)
    \item Dimensionalidad (3)
    \item Topología Klein ($7\pi$)
\end{itemize}
\end{nivelmatematico}

\textbf{LA ECUACIÓN DE ONDA ELECTROMAGNÉTICA:}

De Maxwell se deriva:
\begin{equation}
\nabla^2 \mathbf{E} - \frac{1}{c^2} \frac{\partial^2 \mathbf{E}}{\partial t^2} = 0
\end{equation}

Sustituyendo $c$ Klein:
\begin{equation}
\nabla^2 \mathbf{E} - \frac{1}{\left(3 - \frac{1}{(7\pi)^2}\right)^2 \times 10^{16}} \frac{\partial^2 \mathbf{E}}{\partial t^2} = 0
\end{equation}

Esta es la ecuación que gobierna TODA la luz del universo:
desde las ondas de radio hasta los rayos gamma.


% ============================================================================
\section{Resumen del Capítulo}
% ============================================================================

\begin{tcolorbox}[colback=kleingray, colframe=kleinblue, title=\textbf{Resumen del Capítulo 3}]

\textbf{IDEAS CENTRALES:}

\begin{enumerate}
    \item Las ecuaciones de Maxwell describen toda la luz y el
    electromagnetismo. Dependen de $c$, $\varepsilon_0$, $\mu_0$.

    \item La velocidad de la luz tiene forma Klein:
    \begin{equation}
    c = \left(3 - \frac{1}{(7\pi)^2}\right) \times 10^8 \text{ m/s} \quad \text{[Error: 0.0003\%]}
    \end{equation}

    \item La constante de estructura fina tiene forma Klein:
    \begin{equation}
    \frac{1}{\alpha} = 7^2\pi - 7 - \pi^2 \quad \text{[Error: 0.024\%]}
    \end{equation}

    \item La impedancia del vacío tiene forma Klein:
    \begin{equation}
    Z_0 = 120\pi = 5!\pi \quad \text{[Error: 0.07\%]}
    \end{equation}

    \item Las ecuaciones de Maxwell son consecuencia de la topología Klein.
\end{enumerate}

\vspace{0.3cm}
\textbf{PRÓXIMO CAPÍTULO:}

Si Maxwell viene de Klein, ¿qué pasa con Einstein y la gravedad?

\end{tcolorbox}


% ============================================================================
\section*{Ejercicios}
% ============================================================================

\begin{enumerate}
    \item Verifica que $c = 1/\sqrt{\varepsilon_0\mu_0}$ usando los valores de $\varepsilon_0$ y $\mu_0$.

    \item Calcula qué valor tendría $\alpha$ si la fórmula fuera $1/\alpha = 7^2\pi$ (sin las
    correcciones $-7$ y $-\pi^2$). ¿Qué implicaría para los átomos?

    \item La impedancia $Z_0 = 377$ $\Omega$ es aproximadamente $120\pi$.
    ¿Qué precisión tiene esta aproximación?

    \item Si la velocidad de la luz fuera exactamente $3 \times 10^8$ m/s (sin corrección
    Klein), ¿cuánto diferiría del valor real? ¿Es medible?
\end{enumerate}


% ============================================================================
\section*{Código de Verificación}
% ============================================================================

\begin{lstlisting}[caption={Verificación numérica - Capítulo 3}]
import numpy as np
from numpy import pi

siete_pi = 7 * pi

# Velocidad de la luz
c_obs = 299792458
c_klein = (3 - 1/siete_pi**2) * 1e8
print(f"VELOCIDAD DE LA LUZ:")
print(f"  c observado  = {c_obs} m/s")
print(f"  c Klein      = {c_klein:.0f} m/s")
print(f"  Error = {abs(c_klein - c_obs)/c_obs*100:.6f}%")

# Estructura fina
alpha_inv_obs = 137.035999
alpha_inv_klein = 7**2 * pi - 7 - pi**2
print(f"\nESTRUCTURA FINA:")
print(f"  1/alpha observado = {alpha_inv_obs:.6f}")
print(f"  1/alpha Klein     = {alpha_inv_klein:.6f}")
print(f"  Error = {abs(alpha_inv_klein - alpha_inv_obs)/alpha_inv_obs*100:.6f}%")

# Impedancia
Z_obs = 376.730
Z_klein = 120 * pi
print(f"\nIMPEDANCIA DEL VACIO:")
print(f"  Z0 observado = {Z_obs:.4f} Ohm")
print(f"  Z0 Klein     = {Z_klein:.4f} Ohm")
print(f"  Error = {abs(Z_klein - Z_obs)/Z_obs*100:.4f}%")

# Desglose de la formula de alpha
print(f"\nDESGLOSE DE 1/alpha = 7^2*pi - 7 - pi^2:")
print(f"  7^2*pi = {49*pi:.4f}")
print(f"  -7     = -7.0000")
print(f"  -pi^2  = {-pi**2:.4f}")
print(f"  Total  = {alpha_inv_klein:.4f}")
\end{lstlisting}
